% Options for packages loaded elsewhere
\PassOptionsToPackage{unicode}{hyperref}
\PassOptionsToPackage{hyphens}{url}
\documentclass[
  ukrainian,
]{article}
\usepackage{xcolor}
\usepackage[margin=2.5cm]{geometry}
\usepackage{amsmath,amssymb}
\setcounter{secnumdepth}{-\maxdimen} % remove section numbering
\usepackage{iftex}
\ifPDFTeX
  \usepackage[T1]{fontenc}
  \usepackage[utf8]{inputenc}
  \usepackage{textcomp} % provide euro and other symbols
\else % if luatex or xetex
  \usepackage{unicode-math} % this also loads fontspec
  \defaultfontfeatures{Scale=MatchLowercase}
  \defaultfontfeatures[\rmfamily]{Ligatures=TeX,Scale=1}
\fi
\usepackage{lmodern}
\ifPDFTeX\else
  % xetex/luatex font selection
\fi
% Use upquote if available, for straight quotes in verbatim environments
\IfFileExists{upquote.sty}{\usepackage{upquote}}{}
\IfFileExists{microtype.sty}{% use microtype if available
  \usepackage[]{microtype}
  \UseMicrotypeSet[protrusion]{basicmath} % disable protrusion for tt fonts
}{}
\makeatletter
\@ifundefined{KOMAClassName}{% if non-KOMA class
  \IfFileExists{parskip.sty}{%
    \usepackage{parskip}
  }{% else
    \setlength{\parindent}{0pt}
    \setlength{\parskip}{6pt plus 2pt minus 1pt}}
}{% if KOMA class
  \KOMAoptions{parskip=half}}
\makeatother
\usepackage{longtable,booktabs,array}
\newcounter{none} % for unnumbered tables
\usepackage{calc} % for calculating minipage widths
% Correct order of tables after \paragraph or \subparagraph
\usepackage{etoolbox}
\makeatletter
\patchcmd\longtable{\par}{\if@noskipsec\mbox{}\fi\par}{}{}
\makeatother
% Allow footnotes in longtable head/foot
\IfFileExists{footnotehyper.sty}{\usepackage{footnotehyper}}{\usepackage{footnote}}
\makesavenoteenv{longtable}
\usepackage{graphicx}
\graphicspath{{/Users/kalinicenkovitalijmikolajovic/.codex/worktrees/5c15/Script_for_translate_book/Output_Final/}}
\makeatletter
\newsavebox\pandoc@box
\newcommand*\pandocbounded[1]{% scales image to fit in text height/width
  \sbox\pandoc@box{#1}%
  \Gscale@div\@tempa{\textheight}{\dimexpr\ht\pandoc@box+\dp\pandoc@box\relax}%
  \Gscale@div\@tempb{\linewidth}{\wd\pandoc@box}%
  \ifdim\@tempb\p@<\@tempa\p@\let\@tempa\@tempb\fi% select the smaller of both
  \ifdim\@tempa\p@<\p@\scalebox{\@tempa}{\usebox\pandoc@box}%
  \else\usebox{\pandoc@box}%
  \fi%
}
% Set default figure placement to htbp
\def\fps@figure{htbp}
\makeatother
\ifLuaTeX
\usepackage[bidi=basic,shorthands=off,provide=*]{babel}
\else
\usepackage[bidi=default,shorthands=off,provide=*]{babel}
\fi
\ifLuaTeX
  \usepackage{selnolig} % disable illegal ligatures
\fi
\setlength{\emergencystretch}{3em} % prevent overfull lines
\providecommand{\tightlist}{%
  \setlength{\itemsep}{0pt}\setlength{\parskip}{0pt}}
\usepackage{bookmark}
\IfFileExists{xurl.sty}{\usepackage{xurl}}{} % add URL line breaks if available
\urlstyle{same}
\hypersetup{
  pdflang={uk},
  hidelinks,
  pdfcreator={LaTeX via pandoc}}

\author{}
\date{}

\begin{document}

\section{РОЗВ'ЯЗАННЯ ДИФЕРЕНЦІАЛЬНИХ РІВНЯНЬ У ПОХІДНИХ ПОХІДНИХ ЗА
ДОПОМОГОЮ МЕТОДУ КОНФОРМНИХ
ВІДОБРАЖЕНЬ}\label{ux440ux43eux437ux432ux44fux437ux430ux43dux43dux44f-ux434ux438ux444ux435ux440ux435ux43dux446ux456ux430ux43bux44cux43dux438ux445-ux440ux456ux432ux43dux44fux43dux44c-ux443-ux43fux43eux445ux456ux434ux43dux438ux445-ux43fux43eux445ux456ux434ux43dux438ux445-ux437ux430-ux434ux43eux43fux43eux43cux43eux433ux43eux44e-ux43cux435ux442ux43eux434ux443-ux43aux43eux43dux444ux43eux440ux43cux43dux438ux445-ux432ux456ux434ux43eux431ux440ux430ux436ux435ux43dux44c}

МЕТА ЛЕКЦІЇ: Показати, як деякі двовимірні граничні задачі можна
перетворити на інші, простіші задачі за допомогою конформного
відображення. Припустимо, наприклад, що ми хочемо розв'язати рівняння
Лапласа

\[
\begin{cases} u(0, t) = 0, \\ u_x(1, t) = -u(1, t), \\ u(x, 0) = x, \\ u_t(x, 0) = 0. \end{cases}
\]

з певною граничною умовою в області комплексної форми (площина змінних x
і y). Цю задачу можна перетворити на нову, в якій потрібно знайти
розв'язок рівняння Лапласа

у простішій області (площина змінних і і і).

\pandocbounded{\includegraphics[keepaspectratio]{images/_page_353_Figure_9.jpeg}}

\begin{verbatim}
У конформному відображенні рівняння Лапласа \varphi_{xx}+\varphi_{yy}=0 у координатній площині (x,y) повертається у рівняння Лапласа \varphi_{nn}\dashv -\varphi_{vo}=0 у координатній площині (u,v). Інші рівняння, такі як \varphi_{xx}+2\varphi_{xy}+\varphi_{yy}=0 , не мають такої властивості. Після розв'язання рівняння Лапласа \varphi(u,v) у простій області (
\end{verbatim}

\begin{verbatim}
наприклад, у колі, напівплощині, квадраті) знайдено, достатньо підставити вирази U=U(X, Y), V=V(X, Y) у цей розв'язок, і ми отримаємо розв'язок нашої задачі \phi[u(x, y), v(x, y)] , вираженого через початкові змінні.
\end{verbatim}

У цій лекції ми розглянемо:

\begin{enumerate}
\def\labelenumi{\arabic{enumi})}
\item
  про властивості конформних відображень,
\item
  як їх будувати (для деяких районів),
\item
  як розв'язувати задачі методом конформних відображень.
\end{enumerate}

Однією з головних труднощів, що виникають при розв'язанні задач меж, є
складна форма кордонів. Навіть межі відносно простої форми часто роблять
задачу дуже складною для розв'язання. Звісно, можна спробувати
розв'язати задачу за допомогою методів теорії збурень, але теорія
збурень працює добре лише у випадках, коли межа близька до простої.

Однак існує спосіб розв'язати двовимірне рівняння Лапласа в областях з
комплексною межею. Цей метод базується на конформних відображеннях.

Перш ніж почати розв'язувати задачі цим методом, давайте трохи
ознайомимося з поняттям конформного відображення та комплексних змінних
загалом.

\section{Комплексні функції та конформні
відображення}\label{ux43aux43eux43cux43fux43bux435ux43aux441ux43dux456-ux444ux443ux43dux43aux446ux456ux457-ux442ux430-ux43aux43eux43dux444ux43eux440ux43cux43dux456-ux432ux456ux434ux43eux431ux440ux430ux436ux435ux43dux43dux44f}

У цій лекції ми ознайомимося лише з тими поняттями з теорії функцій
комплексної змінної, які нам знадобляться пізніше. Представимо
комплексне число z = x + iy як точку на площині комплексного xy (рис.
47.1).

Щоб ввести поняття конформного відображення, спершу ознайомимося з
поняттям функції комплексної змінної

.

Тут z --- це комплексна змінна (з деякої області площини z), а w ---
нова комплексна змінна, яка отримає значення відповідно до формули w =
\(f(z)\) у новій комплексній площині w.

Наприклад, комплексна функція, визначена в \emph{першому квадранті}
комплексної z-площини, \(w=z^2\) відображає цю область у всю півплощину
v\textgreater0 у комплексній площині w.

Таблиця 47 2

{\def\LTcaptype{none} % do not increment counter
\begin{longtable}[]{@{}
  >{\raggedright\arraybackslash}p{(\linewidth - 4\tabcolsep) * \real{0.2342}}
  >{\raggedright\arraybackslash}p{(\linewidth - 4\tabcolsep) * \real{0.2252}}
  >{\raggedright\arraybackslash}p{(\linewidth - 4\tabcolsep) * \real{0.5405}}@{}}
\toprule\noalign{}
\begin{minipage}[b]{\linewidth}\raggedright
Площина d
\end{minipage} & \begin{minipage}[b]{\linewidth}\raggedright
\(z \longrightarrow z^2\)
\end{minipage} & \begin{minipage}[b]{\linewidth}\raggedright
Площина S
\end{minipage} \\
\midrule\noalign{}
\endhead
\bottomrule\noalign{}
\endlastfoot
0 & & 0 \\
t1 → t & & -1(2) \\
Вісь x Перший квадрант & & Позитивна дійсна піввісь верхньої
півплощини \\
\end{longtable}
}

Таблиця 47.2 показує відповідність, яка встановлюється шляхом
відображення \(w=z^2\) між деякими точками на площинах z і w.

\pandocbounded{\includegraphics[keepaspectratio]{images/_page_355_Figure_5.jpeg}}

Рис. 47.1. Площина комплексної змінної \(\boldsymbol{z}\) та
некольорових формул бігунів Р И С У Н

\pandocbounded{\includegraphics[keepaspectratio]{images/_page_355_Figure_7.jpeg}}

ОК 47.2. Відображення першого квадранта площини \({\bf z}\) у верхню
півплощину \({\bf w}\) .

Інакше кажучи, якщо ми говоримо про комплексісні функції змінної w =
\(f(z)\), то можна розглянути питання, як криві площини z перетворюються
на криві в площині

Це питання не ставиться в теорії функцій дійсної змінної, з якою читач
уже добре знайомий.

Перш ніж рухатися далі, давайте розглянемо, як \(w=z^2\) складне
відображення типів у еквівалентній дійсній формі. Дійсна форма показує,
як початкові координати (x,y) у площині z перетворюються на нові
координати (u,v) у площині w \(w=z^2\) .

\[
\begin{cases} u(x, 0) = x, \\ u_t(x, 0) = 0. \end{cases}
\]

І давайте окремо прирівняємо реальні та уявні частини в цьому
співвідношенні. У результаті ми отримуємо

\[
\begin{array}{ll} u=x^2-y^2,\\ v=2xy, \end{array}
\]

(действительная форма отображения \(w=z^2\) ).

Эта форма понадобится нам позже. В этом примере она интересна тем, что
показывает, как гиперболы в плоскости z переходят в координатные линии u
= const и v = const в плоскости w.

Теперь мы перейдем к изучению одного частного вида функций комплексного
переменного, которые осуществляют конформные отображения.

\section{Определение конформного
отображения}\label{ux43eux43fux440ux435ux434ux435ux43bux435ux43dux438ux435-ux43aux43eux43dux444ux43eux440ux43cux43dux43eux433ux43e-ux43eux442ux43eux431ux440ux430ux436ux435ux43dux438ux44f}

Отображение w=\(f(z)\) комплексиой плоскости z на комплексную плоскость
w называется конформным в точке \(z_{\bullet}\) плоскости z, если
производная \(f'(z_{\bullet}) \neq 0\) . Отображение \(f(z)\) называется
конформным в области D, если \(f'(z) \neq 0\) в каждой точке области D.

Например, отображение \(f(z) = z^2\) конформно всюду, за исключением
точки z = 0, поскольку \(z'(z) = 2z \neq 0\) для всех \(z \neq 0\) . C
другой стороны, отображение \(e^z\) конформно во всей плоскости z, так
как всюду \(f'(z) = e^z \neq 0\) . Какая же польза от конформных
отображений? Ответ состоит в том, что, решая уравнение Лапласа
\(\phi_{xx} + \phi_{yy} = 0\) внутри какой-то области плоскости
переменных x и y, мы можем рассматривать эту плоскость как плоскость
комплексного переменного z. Рассмотрим комплексное отображение w =
\(f(z)\) плоскости z на плоскость w, в которой введены координаты (u,
v). При отображении w = \(f(z)\) уравнение Лапласа
\(\varphi_{xx} + \varphi_{yy} = 0\) преобразуется в некоторое новое
уравнение с частными производными, зависящее от новых координат и и v.
Замечательным является следующее: если отображение w = \(f(z)\) является
конформным в области, где задано уравнение \(\phi_{xy} + \phi_{uy} = 0\)
, то новое уравнение также будет уравнением Лапласа в координатах и и v.
To есть уравнение \(\phi_{xx} + \phi_{uu} = 0\) перейдет снова в
уравнение \(\varphi_{nn} + \varphi_{nn} = 0\) . Так мы приходим к идее
найти такое конформное отображение, которое переводит область со сложной
границей в область с простой границей (напомним, что отображение
\(w=z^*\) переводит границу первого квадранта плоскости z в
действительную ось плоскости w).

Приведем несколько примеров, в которых будет показано, как можно решать
уравнения с частными производными с помощью конформных отображений.

Уравнение Лапласа в верхней полуплоскости.

Предположим, что нам необходимо решить следующую задачу Дирихле в
верхней полувлоскости (рис. 47.3):

(47.1)

\begin{verbatim}
можно найти стандартным методом разделения переменных. При этом оказывается, что решение   u\left(x,\,t\right)  , представленное в виде ряда
\end{verbatim}

Один із способів розв'язання цієї проблеми --- застосувати перетворення
Фур'є до змінної x, але набагато краще конформно відобразити верхню
половину площини у площину (u, v). Читачеві не буде складно побачити, що
конформне відображення

является суперпозицией простейших колебаний. Пространственная часть
\(X_n\)(x) каждого элементарного колебания является решением
(собственной функцией) некоторой краевой задачи Штурма---Лиувилля.

\begin{verbatim}
Итак, мы научились решать волиовое уравнение   u_{II}=c^2u_{xx}   на всей прямой   -\infty < x < \infty  . Формула Даламбера представляет решение в виде суммы двух бегущих волн, распространяющихся в противоположные стороны. Если же рассматривать это же уравнение на ограниченном промежутке 0 < x < L, то бегущих волн уже не станет, так как они будут взаимодействовать с границами. Вместо них возникают другие волны, которые называют стоячими волнами. Посмотрим, например, что случится, если гитарную струну, закрепленную в точках x=0, x=L, привести в движение. Для ответа на этот вопрос нам необходимо решить следующую смешанную задачу гиперболического типа:
\end{verbatim}

Вона має такі властивості (див. RNS. 47.4).

\begin{enumerate}
\def\labelenumi{\arabic{enumi}.}
\tightlist
\item
  Вона відображає верхню півплощину площини z у смуги
  \(-\infty < u < \infty\) , \(0 < v < \pi\) , у площині w.
\end{enumerate}

\pandocbounded{\includegraphics[keepaspectratio]{images/_page_357_Figure_11.jpeg}}

РИС 47.3 Задача граничного значення у верхній половині площини.

\begin{itemize}
\tightlist
\item
  \begin{enumerate}
  \def\labelenumi{\arabic{enumi}.}
  \setcounter{enumi}{1}
  \tightlist
  \item
    Ребро прямої y=0, -1 \textless{} x \textless{} 1, у площині z (де
    потенціал дорівнює 1), переходить у пряму \(v=\pi,\) -
    \(\infty < u < \infty\) , у площині w. - 3. Два промені
    x\textgreater1, y=0, і x\textless-1, y=0, у площині z проходять
    відповідно у позитивний і негативний
  \end{enumerate}
\end{itemize}

Читач повинен бути впевнений, що може підтвердити ці твердження.

\pandocbounded{\includegraphics[keepaspectratio]{images/_page_358_Figure_2.jpeg}}

Рис. 47.4. Перетворення складної задачі на легку за допомогою
конформного відображення.

Важливість цих властивостей полягає в тому, що вони перетворюють
початкову вадаху (47.1) на дуже просту задачу

(47.2)

\begin{verbatim}
Розв'язок очевидний: \varphi(u, v) = \frac{1}{\pi}v .
\end{verbatim}

\begin{verbatim}
Тепер, щоб знайти розв'язок початкової задачі, потрібно пояснити v через x і y і підставити цей вираз на q(u,v)=v'\pi . Після цього ми отримуємо
\end{verbatim}

Оказывается, в этом случае бегущие волны отражаются от границ таким
образом, что результирующие колебания становятся не бегущими, а
сохраняющими форму в одном положении, т. е.

\pandocbounded{\includegraphics[keepaspectratio]{images/_page_147_Picture_5.jpeg}}

РИС. 20.1. Три типичных стоячих волны \(X(x)\) \(T(t)\).

превращаются в стоячие волны Несколько примеров стоячих волн показано на
рис. 20.1. Если известны профили этих стоячих волн \(X_n(x)\) и характер
колебаний каждой волны \(T_n(t)\) , то решение задачи о колебаниях
гитарной струны представляет собой суперпозицию простейших колебаний
\(X_n(x)\,T_n(t)\) и запи-

сывается в виде

\[
= \ln\left[\frac{z-1}{z+1}\right] + i \arg\left[\frac{z-1}{z+1}\right]
\]

Единственное, что остается сделать нам,--- это подобрать коэффициенты
\(c_n\) так, чтобы при t=0 удовлетворялись начальные условия
\(u\left(x,\ 0\right)=f\left(x\right)\) и
\(u_t\left(x,\ 0\right)=g\left(x\right)\) .

Задачу о колебаниях гитарной струны мы решим методом разделения
переменных.

\section{Решение задачи о колебаниях ограниченной струны методом
разделения
переменных}\label{ux440ux435ux448ux435ux43dux438ux435-ux437ux430ux434ux430ux447ux438-ux43e-ux43aux43eux43bux435ux431ux430ux43dux438ux44fux445-ux43eux433ux440ux430ux43dux438ux447ux435ux43dux43dux43eux439-ux441ux442ux440ux443ux43dux44b-ux43cux435ux442ux43eux434ux43eux43c-ux440ux430ux437ux434ux435ux43bux435ux43dux438ux44f-ux43fux435ux440ux435ux43cux435ux43dux43dux44bux445}

Чтобы решить смешанную задачу

(УЧП)

\[
v = \arg\left[\frac{z-1}{z+1}\right] = \arg\left[\frac{x+iy-1}{x+1+iy}\right] =
= \arg\left[\frac{2y}{x^2+y^2-1}\right]
\varphi(x, y) = \frac{1}{\pi} \arctan\left[\frac{2y}{x^2 + y^2 - 1}\right].
\begin{cases} u(0, t) = 0, \\ u(L, t) = 0, \end{cases}
x^{2} + y^{2} = 1,
\]

якщо потенціал на внутрішньому колі дорівнює одинці, а на зовнішньому
колі дорівнює двом (рис. 47.5). Іншими словами, нам потрібно розв'язати
проблему

(47.3)

Проблема полягає в тому, щоб знайти конформне відображення, яке
перетворює цю область на щось простіше, де нашу проблему легко
розв'язати. У цьому випадку очевидно, що ця проста ділянка має бути
кільцем. Відповідна експозиція Р И С У Н

\pandocbounded{\includegraphics[keepaspectratio]{images/_page_359_Figure_10.jpeg}}

ОК 47.5. Конформне відображення на кільце.

Це вже не так очевидно, але існують довідники, де можна знайти дисплей
із вказаними властивостями. Одним із найзагальних є «Довідник конформних
відображень»

G. Kober 1), яка містить сотні конформних відображень між різними
регіонами. Якщо читач витратить трохи часу і краще ознайомиться з
конформними таблицями відображення, він знайде відповідне конформне
відображення для задачі (47.3). У нашому конкретному випадку це буде
образ виду

(47.4)

де s=-0,146, а t=-6,85. У своїй еквівалентній дійсній формі це буде

(47.5)

\[
\begin{cases} u(x, 0) = f(x), \\ u_t(x, 0) = g(x), \end{cases}
v = \gamma [y(x-t)],
\gamma = 2t/[(x-t)^2 + y^2].
\]

найдем сначала стоячне волны, т. е. решения вида

\[
\begin{aligned} \varphi_{uu} + \varphi_{vv} &= 0, & \text{в кольче,} \\ (\Gamma Y) & & & & & & & \\ (\Gamma Y) & & & & & & \\ (\Gamma Y) & & & & & \\ & & & & & \\ (\Gamma Y) & & & & & \\ & & & & \\ & & & & \\ & & & & \\ & & & & \\ & & & \\ & & & \\ & & & \\ & & & \\ & & & \\ & & & \\ & & \\ & & \\ & & \\ & & \\ & & \\ & & \\ & & \\ & & \\ & & \\ & & \\ & & \\ & & \\ & & \\ & & \\ & & \\ & & \\ & & \\ & & \\ & & \\ & & \\ & & \\ & & \\ & & \\ & & \\ & & \\ & & \\ & & \\ & & \\ & & \\ & & \\ & & \\ & & \\ & & \\ & & \\ & & \\ & & \\ & & \\ & & \\ & & \\ & & \\ & & \\ & & \\ & & \\ & & \\ & & \\ & & \\ & & \\ & & \\ & & \\ & & \\ & & \\ & & \\ & & \\ & & \\ & & \\ & & \\ & & \\ & & \\ & & \\ & & \\ & & \\ & & \\ & & \\ & & \\ & & \\ & & \\ & & \\ & & \\ & & \\ & & \\ & & \\ & & \\ & & \\ & & \\ & & \\ & & \\ & & \\ & & \\ & & \\ & & \\ & & \\ & & \\ & & \\ & & \\ & & \\ & & \\ & & \\ & & \\ & & \\ & & \\ & & \\ & & \\ & & \\ & & \\ & & \\ & & \\ & & \\ & & \\ & & \\ & & \\ & & \\ & & \\ & & \\ & & \\ & & \\ & & \\ & & \\ & & \\ & & \\ & & \\ & & \\ & & \\ & & \\ & & \\ & & \\ & & \\ & & \\ & & \\ & & \\ & & \\ & & \\ & & \\ & & \\ & & \\ & & \\ & & \\ & & \\ & & \\ & & \\ & & \\ & & \\ & & \\ & & \\ & & \\ & & \\ & & \\ & & \\ & & \\ & & \\ & & \\ & & \\ & & \\ & & \\ & & \\ & & \\ & & \\ & & \\ & & \\ & & \\ & & \\ & & \\ & & \\ & & \\ & & \\ & & \\ & & \\ & & \\ & & \\ & & \\ & & \\ & & \\ & & \\ & & \\ & & \\ & & \\ & & \\ & & \\ & & \\ & & \\ & & \\ & & \\ & & \\ & & \\ & & \\ & & \\ & & \\ & & \\ & & \\ & & \\ & & \\ & & \\ & & \\ & & \\ & & \\ & & \\ & & \\ & & \\ & & \\ & & \\ & & \\ & & \\ & & \\ & & \\ & & \\ & & \\ & & \\ & & \\ & & \\ & & \\ & & \\ & & \\ & & \\ & & \\ & & \\ & & \\ & & \\ & & \\ & & \\ & & \\ & & \\ & & \\ & & \\ & & \\ & & \\ & & \\ & & \\ & & \\ & & \\ & & \\ & & \\ & & \\ & & \\ & & \\ & & \\ & & \\ & & \\ & & \\ & & \\ & & \\ & & \\ & & \\ & & \\ & & \\ & & \\ & & \\ & & \\ & & \\ & & \\ & & \\ & & \\ & & \\ & & \\ & & \\ & & \\ & & \\ & & \\ & & \\ & & \\ & & \\ & & \\ & & \\ & & \\ & & \\ & & \\ & & \\ & & \\ & & \\ & & \\ & & \\ & & \\ & & \\ & & \\ & & \\ & & \\ & & \\ & & \\ & & \\ & & \\ & & \\ & & \\ & & \\ & & \\ & & \\ & & \\ & & \\ & & \\ & & \\ & & \\ & & \\ & & \\ & & \\ & & \\ & & \\ & & \\ & & \\ & & \\ & & \\ & & \\ & & \\ & & \\ & & \\ & & \\ & & \\ & & \\ & & \\ & & \\ & & \\ & & \\ & & \\ & & \\ & & \\ & & \\ & & \end{aligned}
\]

Подставив это выражение в уравнение (20.1) и разделив переменные,
получим два обыкновенных дифференциальных уравнения

\[
\varphi(r) = a \ln r + b, \quad r = u^2 + v^2.
\varphi(u, v) = 0.57 \ln(u^2 + v^2) + 1
\]

\begin{verbatim}
где константа   \lambda   — пока что любое действительное число —   \infty < \lambda < \infty   .
\end{verbatim}

\begin{verbatim}
Результаты исследования решений этих двух уравнений при различных значениях   \lambda   приведены на рис 20.2.
\end{verbatim}

\begin{verbatim}
Теперь мы должны *опбросить* те решения, которые либо неограничены при   t \to \infty  , либо становятся тождественно равными нулю, если потребовать выполнения граничных условий u(0, t) = u(L, t) = 0.
\end{verbatim}

\begin{verbatim}
В качестве упражнения читателю предлагается проверить, что нетривиальные и ограниченные решения можно получить только при отрицательных значениях параметра   \lambda  . Таким образом, наша задача— найти такие значения констант A, B, C, D и отрица-
\end{verbatim}

\begin{verbatim}
тельной константы   \lambda = -\beta^2  , чтобы выражение
\end{verbatim}

(20.2)

\[
\varphi(x, y) = 0.57 \ln(u^2 + v^2) + 1,
\]

удовлетворяло граничным условиям. Если это удастся, то мы будем
располагать множеством простейших колебаний струны, каждое из которых
удовлетворяет уравнению с частными произ- Р И С

\pandocbounded{\includegraphics[keepaspectratio]{images/_page_149_Figure_5.jpeg}}

20.2. Построение стоячих воли при раздичных значениях параметра λ.

водными и граничным условиям. Затем останется только найти такую
линейную комбинацию этих колебаний, чтобы сумма удовлетворяла начальным
условиям при t=0.

Подстановка (20.2) в граничные условия \(u(0, t)\) = \(u(L, t)\) = 0
дает

\[
u(0, t) = X(0)T(t) = D\bigl[A\sin(\alpha\beta t) + B\cos(\alpha\beta t)\bigr] = 0 \Rightarrow D = 0,
\]

и

\[
u(L, t) = X(L)T(t) = C\sin(\beta L)\bigl[A\sin(\alpha\beta t) + B\cos(\alpha\beta t)\bigr] = 0 \Rightarrow \sin(\beta L) = 0.
\]

\begin{verbatim}
Другими словами, константа разделения \beta должна удовлетворять уравнению \sin(\beta L) = 0, откуда
\end{verbatim}

\[
\beta_n = \frac{n\pi}{L}, \qquad n = 0, 1, 2, \ldots
\]

\begin{verbatim}
Заметим, что если во втором уравнении (20.3) положить C=0, то получится тривиальное решение X(x)T(t) \equiv 0. Следовательно, мы нашли последовательность элементарных колебаний струны:
\end{verbatim}

\[
u_n(x, t) = X_n(x)T_n(t) = \sin\!\left(\frac{n\pi x}{L}\right)\left[a_n \sin\!\left(\frac{n\pi \alpha t}{L}\right) + b_n \cos\!\left(\frac{n\pi \alpha t}{L}\right)\right].
\]

(20.4) або

\[
u_n(x, t) = R_n \sin\!\left(\frac{n\pi x}{L}\right)\cos\!\left(\frac{n\pi \alpha (t - \delta_n)}{L}\right).
\]

\begin{verbatim}
Тут a_n, b_n, R_n і \delta_n — довільні сталі. Кожне таке елементарне коливання є стоячою хвилею.
\end{verbatim}

\pandocbounded{\includegraphics[keepaspectratio]{images/_page_150_Picture_5.jpeg}}

Рис. 20.3. Стоячі хвилі \(u_n\)(x, t) = \(X_n\)(x)\(T_n\)(t).

В нашей задаче уравнение и граничные условия линейны и однородны,
поэтому произвольная сумма элементарных стоячих волн также является
решением волнового уравнения и удовлетворяет граничным условиям. Мы
должны теперь найти такую суперпозицию стоячих волн, чтобы она
удовлетворяла также и начальным условиям, тогда получим решение нашей
задачи. Подстановка суммы

Зверніть увагу на зображення.

\pandocbounded{\includegraphics[keepaspectratio]{images/_page_362_Figure_1.jpeg}}

\pandocbounded{\includegraphics[keepaspectratio]{images/_page_362_Figure_2.jpeg}}

\begin{enumerate}
\def\labelenumi{\arabic{enumi}.}
\setcounter{enumi}{3}
\tightlist
\item
  Розв'язати змішану задачу Діріхле-Неймана з кутовим порошком за
  допомогою розчину під кутом 45°
\end{enumerate}

\pandocbounded{\includegraphics[keepaspectratio]{images/_page_362_Figure_5.jpeg}}

ПРИМІТКА. Комплексна функція \(w = \ln z = \ln |z| + i \arg(z)\)
відображає промінь \(\theta = c_1\) у площині z у пряму \(v = c_1\) , а
коло \(r = c_2\) у пряму \(u = \ln c_2\) у площині w (див. рисунки).

\pandocbounded{\includegraphics[keepaspectratio]{images/_page_362_Figure_7.jpeg}}

\section{Джерела ()}\label{ux434ux436ux435ux440ux435ux43bux430}

\section{Головне}\label{ux433ux43eux43bux43eux432ux43dux435}

\begin{itemize}
\tightlist
\item
  Бідадзе А. В. Рівняння математичної фізики.--- Москва: Наука, 1976.
\item
  Володимиров В. С. Рівняння математичної фізики.--- Москва: Наука,
  1971. Кошляков Н. С., Глінер Е. Б., Смірнов М. М. Рівняння в зокрема
  математичній фізиці.--- Москва: Вища школа, 1970.
\item
  \begin{enumerate}
  \def\labelenumi{\arabic{enumi}.}
  \setcounter{enumi}{3}
  \tightlist
  \item
    Самарський, А. А. Вступ до чисельних методів. Москва, Наукове
    видавництво, 1982.
  \end{enumerate}
\item
  \begin{enumerate}
  \def\labelenumi{\arabic{enumi}.}
  \setcounter{enumi}{4}
  \tightlist
  \item
    Смирнов В. І. Курс, виший математіки, т. 2.--- М.: Наука, 1967.
  \end{enumerate}
\item
  \begin{enumerate}
  \def\labelenumi{\arabic{enumi}.}
  \setcounter{enumi}{5}
  \tightlist
  \item
    Тихонов А. Н., Самарський А. А. Рівняння математичної фізики.
    Москва, Наукове видання, 1972.
  \end{enumerate}
\end{itemize}

\section{Додатково}\label{ux434ux43eux434ux430ux442ux43aux43eux432ux43e}

\begin{itemize}
\tightlist
\item
  \begin{enumerate}
  \def\labelenumi{\arabic{enumi}.}
  \setcounter{enumi}{6}
  \tightlist
  \item
    Задача Адамара Дж. Коші для лінійних диференціальних рівнянь у
    похідних похідних гіперболічного типу. Москва, Наукове видання,
    1978.
  \end{enumerate}
\item
  Арсенін В. Я. Методи математичної фізики та спеціальні функції.
  Москва; Наука, 1974.
\item
  \begin{enumerate}
  \def\labelenumi{\arabic{enumi}.}
  \setcounter{enumi}{8}
  \tightlist
  \item
    Арфкен Г. Математичні методи у фізиці.--- Москва: Атомиздат, 1970.
  \end{enumerate}
\item
  \begin{enumerate}
  \def\labelenumi{\arabic{enumi}.}
  \setcounter{enumi}{9}
  \tightlist
  \item
    Берс Л., Джон Ф., Шехтер Б. Диференціальні рівняння з частинними
    похідними. Москва, видавництво Мир, 1966.
  \end{enumerate}
\item
  \begin{enumerate}
  \def\labelenumi{\arabic{enumi}.}
  \setcounter{enumi}{10}
  \tightlist
  \item
    Бінадзе, А. В. Крайові задачі для еліптичних рівнянь другого
    порядку.--- Москва: Наука, 1966.
  \end{enumerate}
\item
  \begin{enumerate}
  \def\labelenumi{\arabic{enumi}.}
  \setcounter{enumi}{11}
  \tightlist
  \item
    Бідадзе А. В. Деякі класи диференціальних рівнянь у похідних
    похідних. Москва, Наукове видання, 1981.
  \end{enumerate}
\item
  \begin{enumerate}
  \def\labelenumi{\arabic{enumi}.}
  \setcounter{enumi}{12}
  \tightlist
  \item
    Біцадзе А. В., Калініченко Д. К. Збірка задач з рівнянь математичної
    фізики. Москва, Наука, 1977.
  \end{enumerate}
\item
  \begin{enumerate}
  \def\labelenumi{\arabic{enumi}.}
  \setcounter{enumi}{13}
  \tightlist
  \item
    Будак Б. М., Самарський А. А., Тихонов А. Н. Сборник задач по
    математичній фізиці {[}Збірка задач математичної фізики{]}. Москва,
    Наукове видання, 1972.
  \end{enumerate}
\item
  Вазов В., Форсайт Дж. Методи різниці для розв'язання диференціальних
  рівнянь у похідних похідних.--- Москва: Іл, 1963
\item
  Нові методи розв'язання еліптичних рівнянь {[}Нові методи розв'язання
  еліптичних рівнянь{]}. Москва, Гостехіддат, 1948. Векуа, І. Н.
  Узагальнені аналітичні функції. Москва, Фізматгиз, 1959.
\item
  Володимиров В. С. Узагальнені функції в математичній фізиці, Москва,
  Наука, 1976.
\item
  \begin{enumerate}
  \def\labelenumi{\arabic{enumi}.}
  \setcounter{enumi}{18}
  \tightlist
  \item
    Володимиров В. С. та ін. Збірка задач з рівнянь математичної фізики.
    --- Москва: Наука, 1974.
  \end{enumerate}
\item
  Годунов С. К. Рівняння математичної фізики.--- Москва: Наука, 1971.
  Годунов С. К., Рябенський В. С. Схеми різниці.--- Москва: Наука, 1977.
\item
  Гордін: задача Л. Коші для гіперболічних рівнянь.--- Москва: IL, 1961.
  Джеффріс Г., Спірлс Б. Методи математичної фізики. т. 1-3, --- Москва:
  Мир,
\item
  \begin{enumerate}
  \def\labelenumi{\arabic{enumi}.}
  \setcounter{enumi}{1969}
  \tightlist
  \item
  \end{enumerate}
\item
  \begin{enumerate}
  \def\labelenumi{\arabic{enumi}.}
  \setcounter{enumi}{23}
  \tightlist
  \item
    Зоммерфельд А. Рівняння в часткових похідних математичної фізики.---
    Москва: Іл, 1950. 25. Кализкін, Н. Н. Численільна методика.--- М.:
    Наука, 1978.
  \end{enumerate}
\item
  25 Курант Р., Гілберт Д. Методи математичної фізики. т. 1-2, - Москва:
  Гостехиздат, 1951.
\item
  \begin{enumerate}
  \def\labelenumi{\arabic{enumi}.}
  \setcounter{enumi}{26}
  \tightlist
  \item
    Курант Р. Рівняння з частковими виробництвами. Москва, Видавництво
    Мир, 1964.
  \end{enumerate}
\item
  \begin{enumerate}
  \def\labelenumi{\arabic{enumi}.}
  \setcounter{enumi}{27}
  \tightlist
  \item
    Лаврентьєв М. А., Шабат Б. В. Методи теорії функцій комплексної
    змінної. - Москва: Наука, 1965.
  \end{enumerate}
\item
  \begin{enumerate}
  \def\labelenumi{\arabic{enumi}.}
  \setcounter{enumi}{28}
  \tightlist
  \item
    Ладиженська О. А. Прикордонні задачі математичної фізики. Москва:
    Наука,
  \end{enumerate}
\item
  Ладиженська О. А. Змішана задача для гіперболічних рівнянь.--- Москва:
  Фізматгиз, 1953.
\end{itemize}

Список був складений редактором перекладу.

\begin{itemize}
\tightlist
\item
  Ладиженська О. А., Солоняков В. А., Уралцева Н. Н. Лінійні та
  кварілінні рівняння параболічного типу.--- Москва: Наука, 1967.
\item
  \begin{enumerate}
  \def\labelenumi{\arabic{enumi}.}
  \setcounter{enumi}{31}
  \tightlist
  \item
    Ладиженська О. А., Уралцева Н. Н. Лінійні та квазілінійні рівняння
    еліптичного типу. - Москва: Наука, 1964, 1973.
  \end{enumerate}
\item
  Лебедєв Н. Н., Скальська М. П., Уфляйнд Ю. С. Сборник задах по
  математичній фізиці.--- Москва: Гостехиддат, 1965.
\item
  \begin{enumerate}
  \def\labelenumi{\arabic{enumi}.}
  \setcounter{enumi}{33}
  \tightlist
  \item
    Лайонс, Дж.--- Л. Деякі методи розв'язання задач з нелінійними
    краями.--- М.З. Мір. 1972.
  \end{enumerate}
\item
  \begin{enumerate}
  \def\labelenumi{\arabic{enumi}.}
  \setcounter{enumi}{34}
  \tightlist
  \item
    Лайонс Дж.---Л., Maggenes E. Neobreed Grancinum Problems та їх
    застосування. Москва, Видавництво Мир, 1971. 36. Марчук Г. І. Методи
    обчислювальної математики. Москва, Наукове видання, 1980.
  \end{enumerate}
\item
  Мізохата С. Теорія рівнянь з приватними виробниками. --- Москва: Мир,
  1977.
\item
  \begin{enumerate}
  \def\labelenumi{\arabic{enumi}.}
  \setcounter{enumi}{37}
  \tightlist
  \item
    Міранда К. Рівняння з частковими похідними еліптичного типу.---
    Москва: Іл, 1957.
  \end{enumerate}
\item
  Михайлов, В. П. Диференціальні рівняння в похідних з похідними. ---
  Москва: Наука, 1976.
\item
  Міклін С. Г. Курс математичної фізики.--- Москва: Наука, 1968. Міхлін,
  С. Г. Лінійні диференціальні рівняння. --- Москва; Вишня школа, 1977.
\item
  Морс П. М., Фецибах X, Методи теоретичної фізики. т. 1-2.--- Москва:
  IL, 1958.
\item
  \begin{enumerate}
  \def\labelenumi{\arabic{enumi}.}
  \setcounter{enumi}{42}
  \tightlist
  \item
    Меттьюз Дж., Вокер Р. Математичні методи фізики.--- Москва:
    Атомиздат, 1972.
  \end{enumerate}
\item
  Лекції Нагумо М. з сучасної теорії рівнянь у приватних виробництвах
  Москва, видавництво Мвр, 1967. 45. Обен, Дж.-П. Приблизне розв'язання
  еліптичних граничних задач.--- Москва:
\item
  Мір, 1977. 46. Петровський І. Г. Лекції з рівнянь із частковими
  похідними. --- Москва: Фізматгиз, 1961.
\item
  \begin{enumerate}
  \def\labelenumi{\arabic{enumi}.}
  \setcounter{enumi}{46}
  \tightlist
  \item
    Положий Г. Н. Еквалізація математичної фізики.--- Москва: Вища
    школа.
  \end{enumerate}
\item
  \begin{enumerate}
  \def\labelenumi{\arabic{enumi}.}
  \setcounter{enumi}{47}
  \tightlist
  \item
    Рід М., Саймон Б. Методи сучасної математичної фізики.--- Москва:
    Мир,
  \end{enumerate}
\item
  т. 1-1977, т. 2-1978, т. 3, 4-1982. Ріхамаєр Р. Принципи сучасної
  математичної фізики.--- Москва: Мир,
\item
  \begin{enumerate}
  \def\labelenumi{\arabic{enumi}.}
  \setcounter{enumi}{49}
  \tightlist
  \item
    Ріхтмаєр Р., Мортон К. Методи різниці для розв'язання задач з
    граничними значеннями.--- Москва:
  \end{enumerate}
\item
  Мір, 1972,
\item
  Рождественський Б. Л., Япенко, Н. Н. Системи квазіполіцейських рівнянь
  та їх застосування в газовій динаміці.--- Москва: Наука, 1978.
\item
  52, Самарський А. А. Теорія різницьких схем. Москва, Наука, 1977.
\item
  \begin{enumerate}
  \def\labelenumi{\arabic{enumi}.}
  \setcounter{enumi}{52}
  \tightlist
  \item
    Самарський А. А., Андрєєв В. Б. Різниці для еліптичних рівнянь---
    Москва, Наука, 1976.
  \end{enumerate}
\item
  \begin{enumerate}
  \def\labelenumi{\arabic{enumi}.}
  \setcounter{enumi}{53}
  \tightlist
  \item
    Смірнов В. І. Курс вищої математики.--- Москва: Наука, 1974.
  \end{enumerate}
\item
  \begin{enumerate}
  \def\labelenumi{\arabic{enumi}.}
  \setcounter{enumi}{54}
  \tightlist
  \item
    Смірнов М. М. Дегенеративні еліптичні та гіперболічні рівняння.--
    Москва: Наука, 1966.
  \end{enumerate}
\item
  56, Смірнов М. М. Диференціальні рівняння у похідних другого порядку,
  -- Мінськ: Видавництво БДУ, 1974.
\item
  \begin{enumerate}
  \def\labelenumi{\arabic{enumi}.}
  \setcounter{enumi}{56}
  \tightlist
  \item
    Смірнов М. М. Проблеми рівнянь математичної фізики. Москва, Наукове
    видавництво, 1965.
  \end{enumerate}
\item
  \begin{enumerate}
  \def\labelenumi{\arabic{enumi}.}
  \setcounter{enumi}{57}
  \tightlist
  \item
    Соболєв С. Л. Пекаки primeneniya funktsional'nogo analiza v
    matematicheskoi fizike {[}Застосування функціонального аналізу в
    математичній фізиці{]}. Ленінград: Видавництво ЛГУ, 1950,
    Новосибірськ: СБ АН РСРСР, 1962.
  \end{enumerate}
\item
  Соболєв, С. Л. Рівняння математичної фізики.--- М · Наука, 1966.
\item
  Траптер К. Дж. Інтегральні перетворення у математичній фізиці.---
  Москва: Гостехиддат, 1957.
\item
  \begin{enumerate}
  \def\labelenumi{\arabic{enumi}.}
  \setcounter{enumi}{60}
  \tightlist
  \item
    Трев Ф. Лекції з теорії диференціальних рівнянь у похідних
    похідних.--- Москва: Мир, 1965.
  \end{enumerate}
\item
  Трікомі Ф. Лекції з диференціальних рівнянь у похідних похідних.
  Москва, IL Publ., 1957. Фрідман А. Рівняння з частковими похідними
  параболічного типу. --- Москва: Мир, 1968.
\item
  Германдер Л. Лінійні диференціальні оператори з частковими
  похідними.--- Москва: Мир. 1965.
\item
  \begin{enumerate}
  \def\labelenumi{\arabic{enumi}.}
  \setcounter{enumi}{64}
  \tightlist
  \item
    Шварц Л. Математичні методи фізичних наук.--- Москва: Мир, 1965.
  \end{enumerate}
\item
  \begin{enumerate}
  \def\labelenumi{\arabic{enumi}.}
  \setcounter{enumi}{65}
  \tightlist
  \item
    Шилов Г. Е. Математичний аналіз. Спеціальний курс. Москва,
    Фізмат-гіз видавництво, 1960.
  \end{enumerate}
\item
  \begin{enumerate}
  \def\labelenumi{\arabic{enumi}.}
  \setcounter{enumi}{66}
  \tightlist
  \item
    Щилов, Г. Е. Математичний аналіз. Другий спеціальний курс.---
    Москва: Наука, 1965.
  \end{enumerate}
\item
  \begin{enumerate}
  \def\labelenumi{\arabic{enumi}.}
  \setcounter{enumi}{67}
  \tightlist
  \item
    Яненко Н. Н. Метод дробових щагів для розв'язання багатовимірних
    задач математичної фізики.--- Новосибірськ: Наука, 1967.
  \end{enumerate}
\end{itemize}

\section{ІНДЕКС
ІМЕНІ}\label{ux456ux43dux434ux435ux43aux441-ux456ux43cux435ux43dux456}

Бернуллі 329, Бессель 106, Буффон 322, Ніколсон 296, Ньютон 296

Гамільтон 335, 336 Гаусс 286 Гельмгольц 223 Грін 10 Гук 147 Гюйгенс 177,
181 Пуассон 180, 235

Д'Аламбер 126--129, Діріхле 85, 240, Дюамель 107, 108, 110--112, Рітц
338, 340, Робін 240, Родюг 268

** Seidel 286

Сімсон 320 Стокс 8, 180

Кіз 305, 306 Крхгоф 124

Кірхгоф 124 Кранк 296 Крістоффель 363 Ферма 335 Фур'є 24, 25, 36, 41,
74, 84, 91

Гакслі 304, Ханкел 105, 106, Хевісайд 82, 115, 116, Ходжкін 304-306

Лагранж 329 Ламм 307 Лаплас 8, 50, 98-107 Лежандр 235 Лельбніц 180
Лібман 285 Ліувіль 51, 58, 66, 148

Шварц 363 Шредінгер 8 Штурм 51, 58, 66, 148

Ейлер 84, 215, 329

\textbf{М}аксвелл 8 Меллін 105, 106

10sr (Молодий) 124

Navier 8 Neumann 240, 286

Jacobi 286

\section{ІНДЕКС
ПРЕДМЕТІВ}\label{ux456ux43dux434ux435ux43aux441-ux43fux440ux435ux434ux43cux435ux442ux456ux432}

Матрична алгебра 215, 216 Аналітичне розв'язання 301 Наближення сплайну
10

Позачасовий час 165

Калькулюс варіацій 329 Методи варіацій 10 Зовнішня похідна 79 Хвилі, що
рухаються, 149

-- квартира 177 -- стоячи 222

сферичний 178

--- ударні 210 циліндричні 177, 182

Гармонік 84, 154 Сферичний 265 градієнів 22 Граничні умови 17, 22, 23,
26

--- лінійний 26

--- однорідні 36, 46 --- гетерогенні 26, 49, 50

Дельта-функція 276 Діагоналізація матриці 216 Задача дифузії 113, 163
Дифузійний потік 114 --- член 113 Дифузія 113

Задача ** 147 -- Діріхле 242, 243, 246

\begin{itemize}
\item
  Діріхле - Нойман 364 - Кірхгоф 124
\item
  Коуу 69, 95, 116, 126, 207
\end{itemize}

--- мінімізація 10 --- Нойман 243---245

Ньютон 52, 121

--- на брахістохроні 329 Штурму --- Ліувіль 51, 58, 59, 66,

Закон збереження енергії 31, 114, 205

--- Теплопередача Ньютона 8 246

--- Фур'є 24, 25 Звукові хвилі 124 Момент згинання 157 Інтеграл Фур'є 87
Інтегральні перетворення 9, 73.76. Рівняння 10

Класифікація UCP 308 Крутильні осциляції 123 Поздовжні 123 Конвективні
хвилі 116 Конвективний транспорт 113 --- Член 113 Конвекція 113
Скінченні різниці 282 Канонічні координати 127, 132 --- Полярні 233,
235-238 Сферичні 233, 235, 236, 239, 266 Циліндричні 233, 235, 236, 239
--- еліптичний 239 Неперервний набір 96 Косинусне перетворення 73 ---
скінченне 185, 186 --- Обернене 75, 185, 186 Коефіцієнт дифузії 119 ---
теплопередача 23 Задача граничного значення 240, 241

Лапласіан 233 Лінійність 11

--- --- Еліптичний тренажер 241

Метод дедукції 319 -- Лібман 285, 286 -- Монте-Карло 316 -- Рітц 338,
340 -- Сімпсон 320 спуск 182 -- Фур'є 36 Режим коливання 154 Модуль Юнга
124, 145

Початковий зсув 231 Початкові умови 18

Оператор Лапласа 233 Ортогональність 41, 84, 90

\begin{itemize}
\tightlist
\item
  Ходжкіна - Гакслі 304
\end{itemize}

Власні вектори 67, 215, 217

Випадкові номери 319

Змішане завдання 19

--- enachenia 54, 215, 216 Відгук 96 функція 10, 65 Конформне
відображення 253, 356 Крістоффель-Шварц 363 --- частота 162 Спектральна
функція 76, 77 --- неперервна 77 --- частота 85 Stock 274 Параметрична
ідентифікація 304 Змінна залежна 8 Стоячі хвилі 148, 149 Суперпозиція
97, 193, 223 --- незалежна 8 Схема підрахунку тікання 289 Щільність
матеріалу 30 --- явна 289 Многочлен Лежандра 265 Кранк--Ніколсон 296,
299, 328 Порядок рівняння 10 Потенціал Pewtop 266 з круговою симетрією
240. сферично-симетрична 240 циліндрично симетрична 270 Температурна
характеристика 69 Теорема Діріхле 85 Правило Лейбніца 180--Стокс 180
Перетворення залежних змінних Теорія збурень 10 Тепловий імпульс 96
інтеграл 9, 73, 76 Теплоємність 30 --- координати 9 --- Лаплас 50, 75,
98---107, 115, 116, теплопровідність 25, 30, 33 126 Точкове джерело 274
--- обернене 75 --- Меллін 105, 106 --- Ханкель 76, 105, 106 ---
обернене 76. Питома теплоємність 34 --- Фур'є 74, 75, 83, 85, 89, 91
Вирівнювання Бесселя 106, 224, 239 --- хвиля 8, 142, 169, 234 --- ---
скінченне 185 --- нелінійне 304 --- пряме 88, 116, 126 --- обернене 88,
74, 75 гармонічні коливання 223 --- Гельмгольц 223, 224, 235 Принцип
Гамільтона 335, 336 --- Гюйгенс 177, 181 --- Дюамель 107, 108, 110, 112
Рух Ньютона 8 дифузія 18, 197 - Суперлосоцитності 191, 196, 250
конвектива 19, 63, 113, 119, 120 ** - Ферма 335** - Ламм 307 - Лаплас 8,
170, 235, 258, 285, 363 - Лежандр 239, 268 Змінне розділення 9, 191
лінійне 11 - гіперболічне 11, 12, 121, 169 Розклад власних функцій 10,
66 - параболічне 11, 12, 14, 169 - еліптичне 11, 12, 169 - Максвелл 8 -
Нав'є - Стокс 8 Згортка 90, 93 неліцейних {[}{]} гетерогенних 11, 19,
65. --- скінченна 102 Гравітаційна сила 122 гомогенна P Пуассон 235. --
--- ранкового перетворення 122 синусоїдального перетворення 75 постпост
285 скінченне 185, 186, 193 зворотне 75, 185, 186 реакції з нелінійною
дифузією Система телеграфних рівнянь 124 --- з прямокутними
коефіцієнтами 60.

170

з частковими похідними 7

Телеграф 8, 123, 125, 215

Зберігання тепла 16

{\def\LTcaptype{none} % do not increment counter
\begin{longtable}[]{@{}llllllll@{}}
\toprule\noalign{}
\_ & tep & Lonre & Доступність води & 8, & 16, & 30, & 169. \\
\midrule\noalign{}
\endhead
\bottomrule\noalign{}
\endlastfoot
& & 234. & & • & • & • & , \\
\end{longtable}
}

\begin{itemize}
\tightlist
\item
  характеристики 202, 208
\end{itemize}

--- Шредінгер, 8 років

--- Ейлер 215, 256, 268

--- Ейлер-Лагранж 329, 330, \textbf{33}2

Стан Діріхле 240

\begin{verbatim}
- Neumann 240 - Robin 240
\end{verbatim}

Сингулярний диференціал 268

Гріна 10, 96, 196, 273, 276

д'Аламбер 126-132

\begin{itemize}
\tightlist
\item
  додатково до ймовірнісного інтеграла 80
\end{itemize}

Неточника 96, 276

-- неперіодичні 77 -- періодичні 77 -- Хевісайд 82, 115, 116, 365.
Фур'є-образ 88, 89, 91, 93

Формула д'Аламбера 180, 183

--- Лейбніц 184

Джерела 99, 100

\begin{itemize}
\tightlist
\item
  Пуассон 180
\end{itemize}

--- --- інтеграл 248, 249, 251

--- Родріга 268

--- Ейлер 84, 85 Фундаментальне рішення 37 Функціональність 329

Функція Бесселя 224, 225

Чисельні методи 9 Кількість персемних 11

Рівнопотенціальні лінії 277

Ядро перетворення 73 -- перетворення Пуассона 252

\section{ЗМІСТ}\label{ux437ux43cux456ux441ux442}

{\def\LTcaptype{none} % do not increment counter
\begin{longtable}[]{@{}
  >{\raggedright\arraybackslash}p{(\linewidth - 4\tabcolsep) * \real{0.1856}}
  >{\raggedright\arraybackslash}p{(\linewidth - 4\tabcolsep) * \real{0.1237}}
  >{\raggedright\arraybackslash}p{(\linewidth - 4\tabcolsep) * \real{0.6907}}@{}}
\toprule\noalign{}
\begin{minipage}[b]{\linewidth}\raggedright
З редагування
\end{minipage} & \begin{minipage}[b]{\linewidth}\raggedright
Гора
\end{minipage} & \begin{minipage}[b]{\linewidth}\raggedright
Переклад
\end{minipage} \\
\midrule\noalign{}
\endhead
\bottomrule\noalign{}
\endlastfoot
Передмова & viye & ****** \\
& & \\
Частина 1. & Введено & \\
Лекція & 1. & Вступ до теорії диференціальних рівнянь з похідними
похідними \\
Частина 2. & Діф & Завдання з синтезу \\
Лекція & 2. & Задачі типу дифузії (параболічні рівняння) \\
Лекція & 3. & Граничні умови в задачах дифузії \\
Лекція & 4. & Виведення рівняння теплопровідності , \\
Лекція & 5. & Розщеплення змінних \\
Лекція & 6. & Перетворення неоднорідних граничних умов у одно- \\
& & Родичі \\
Лекція & 7. & Розв'язання складніших задач методом розділення змін- \\
Тторене \textasciitilde{} & o & Перетворити комплексні рівняння у просту
форму \\
Лекція-лекція & o. & Перетворення комплексних рівнянь у простий
вигляд \\
Мексика & σ, & Власні функції , \\
Pection & 10 & Інтегральні перетворення (перетворення синусів і
козіпус- \\
Боулз & 10. & ) \\
Лекція & 11. & Перетворення рядів і четвірок \\
Лекіня & 12 & Перетворення Фур'є та його застосування до розв'язку \\
* 10HM 11 /1 & , & Рівняння з похідними похідними \\
Лекція & 13. & Перетворення Лапласа \\
Лекція & 14. & Принцип Дюамеля \\
Лекчич & 15. & Принцип Дюамеля \\
& & \\
Частина 3. & Guipe & Проблеми \\
Лекпея & 16. & Одномірне хвильове рівняння (гіперболічні еквалізації) \\
& & \\
Лекція & 17. & Формула Д'Аламбера \\
Лекція & 18. & Формула Д'Аламбера (продовження) \\
Лекція & 19. & Рівняння води та граничні умови \\
Лекція & 20. & Вібрації обмеженої струни (стоячі хвилі) 14 \\
Лекція & 21. & Осциляції пучка (рівняння часткового виробника \\
& & Четвертий порядок) \\
& & Перехід на безрозмірні змінні \\
Лекція & 23. & Класифікація диференціальних рівнянь у похідних \\
& & (Канонічна форма гіперболічного рівняння) 10 \\
Лекція & 24. & Хвильове рівняння у вільному просторі (двовимірне \\
\_ & \textasciitilde{} - & та тривимірні задачі) \\
Лекція & 25. & Скінченні перетворення Фур'є (синус і косинус до - \\
& 00 & Освіта) \\
& & Принцип суперпозиції --- основи теорії динамічних систем 19 \\
& & Рівняння першого порядку (характеристичний метод) 19 \\
Лекція & 28. & Нединамічні рівняння першого порядку (закони
збереження \\
t. & 00 & Ньяя) , \\
Лекція & 29. & Системи диференціальних рівнянь у похідних 2 \\
Розділ & & Мембранні осциляції (покращення хвиль у Polar 2000
Godinates) \\
& & БУУУУММаалака аа аа аа аа аа \\
\end{longtable}
}

{\def\LTcaptype{none} % do not increment counter
\begin{longtable}[]{@{}
  >{\raggedright\arraybackslash}p{(\linewidth - 2\tabcolsep) * \real{0.7658}}
  >{\raggedright\arraybackslash}p{(\linewidth - 2\tabcolsep) * \real{0.2342}}@{}}
\toprule\noalign{}
\begin{minipage}[b]{\linewidth}\raggedright
Частина 4. Еліптичні задачі
\end{minipage} & \begin{minipage}[b]{\linewidth}\raggedright
233
\end{minipage} \\
\midrule\noalign{}
\endhead
\bottomrule\noalign{}
\endlastfoot
Лекція 31. Лапласнан (інтутивна олісанація) & 233240248255 \\
гармоніки) & 264273 \\
Частина 5. Чисельні та наближені методи & 282 \\
Лекція 37. Чисельні розв'язки (еліптичні задачі) & 282289 \\
Lsson) & 296301 \\
Рівняння) & 308316 \\
Монте-Карло & 322329 \\
Лекція 45. Варіаційні методи розв'язання диференціальних рівнянь у
похідних похідних & 338 \\
Лекція 46. Розв'язок рівнянь часткових похідних теоретиків збурень &
346 \\
Лекція 47. Розв'язання диференціальних рівнянь у похідних похідних за
методом конформних відображень & 355 \\
Додаток 1 & 365 \\
Додаток 2 & 371 \\
Spinning 3 & 373 \\
Джерела & 375 \\
Індекс назв & 378 \\
Індекс предметів & 379 \\
\end{longtable}
}

\section{Дорогий
читачу!}\label{ux434ux43eux440ux43eux433ux438ux439-ux447ux438ux442ux430ux447ux443}

Будь ласка, надсилайте свої коментарі щодо змісту приладів, його
конструкції, якості перекладу та іншого за адресою: 129820, Москва,
I-110, GSP. 1-й Рижський провулок, 2, видавництво «Мир».

\section{Фарлоу
Стейпле}\label{ux444ux430ux440ux43bux43eux443-ux441ux442ux435ux439ux43fux43bux435}

\section{ДИФЕРЕНЦІАЛЬНІ РІВНЯННЯ В ПОХІДНИХ ПОХІДНИХ ДЛЯ НАУКОВЦІВ І
ІНЖЕНЕРІВ}\label{ux434ux438ux444ux435ux440ux435ux43dux446ux456ux430ux43bux44cux43dux456-ux440ux456ux432ux43dux44fux43dux43dux44f-ux432-ux43fux43eux445ux456ux434ux43dux438ux445-ux43fux43eux445ux456ux434ux43dux438ux445-ux434ux43bux44f-ux43dux430ux443ux43aux43eux432ux446ux456ux432-ux456-ux456ux43dux436ux435ux43dux435ux440ux456ux432}

Науковий редактор М. В. Хатунцева-молодший науковий редактор Е. І.
Качулина Художник В. Е. Карпов Художній редактор В. І. Шаповалов
Технічний редактор Т. А. Максимова Коректор Т. І. Стіфсева

\section{UB No 5289}\label{ub-no-5289}

Набор 27.12.84. Подано до публікації 01.07.85. Формат \(60\times90^{1}\)
є. Папір книги. журналу. Печатка висока. Літературний Гарінтура. Том L2
--- бум. n, Usl. Печ. с. 24. Стан. Cr.-Ott. 24. Уч.-изд. л. 19.69, Ізд.
No 1/3807. Тираж 5 000 примірників. Замовлення L6 241. Ціна 1 ранд. 70
тис.

MIR 129820, GSP, Москва, I-110, 1-й Рижський провулок, 2

Надруковано в Ленінградській друкарні No 6 Ордена Червоного Прапора
Праці Ленінградської асоціації «Технічний король» імені Євга. Соколова
Союзполіграф при Державному комітеті СРСР з видавничої справи,
друкарства та книжкової торгівлі, 193144, Ленінград, вулиця Моісеєнко,
10, з матрицями Ордена Жовтневого Регулюта та Ордена Червоного Прапора
Праці МПО «Перша зразкова друкарня» імені А. А. Жданога
Союзполіграфірому при Державному комітеті СССГ з видавничої, друкарської
та книжкової торгівлі, 113054, Москам, Валовія, 28

\end{document}
